\section{Pure \texorpdfstring{$2{:}1$}{2:1} parametric forcing:
\texorpdfstring{$f=0$}{f=0}}
\label{sec:all-forced}

\subsection{Why the reduced additive term is exactly absent}

Use the weakly nonlinear scaling
\begin{equation}
 u=\epsilon u_1+\epsilon^2u_2+\epsilon^3u_3+O(\epsilon^4),
 \qquad T=\epsilon^2t.
 \label{eq:pure21-scaling}
\end{equation}
Each separated input is now restricted to one of two choices:
\begin{center}
\begin{tabular}{@{}lll@{}}
\toprule
Input & Present & Absent\\
\midrule
Temporal &
$\epsilon^2(P\ee^{2\ii\omega_0t}+\overline P\ee^{-2\ii\omega_0t})$ &
$0$\\
Spatial &
$\epsilon^2(S\ee^{2\ii k_0x}+\overline S\ee^{-2\ii k_0x})$ &
$0$\\
\bottomrule
\end{tabular}
\end{center}
Thus $n=2$ when the temporal input is present and $m=2$ when the spatial
input is present; \emph{absent} may equivalently be denoted by an infinite
order.  The label $2{:}1$ is a resonance ratio, not an amplitude ratio: the
forcing frequency or wavenumber is twice the corresponding critical carrier,
\begin{equation}
 \omega_F=2\omega_0,\qquad k_F=2k_0.
 \label{eq:pure21-definition}
\end{equation}
The critical first-order field is
\begin{equation}
 u_1=A(T)\Phi\ee^{\ii(\omega_0t+k_0x)}
 +B(T)\Phi\ee^{\ii(\omega_0t-k_0x)}+\mathrm{c.c.}
 \label{eq:pure21-critical-field}
\end{equation}
The two forcing responses in $u_2$ are
\begin{align}
 p_2&=(2\ii\omega_0I-M_0(0))^{-1}P,\notag\\
 s_2^+&=-M_0(2k_0)^{-1}S,\qquad
 s_2^-=-M_0(-2k_0)^{-1}\overline S .
 \label{eq:pure21-responses}
\end{align}
At order $\epsilon^3$, the quadratic term
$2\cQ(u_1,u_2)$ contains
\begin{align}
 (2\omega_0,0)+(-\omega_0,k_0)&=(\omega_0,k_0),\notag\\
 (2\omega_0,0)+(-\omega_0,-k_0)&=(\omega_0,-k_0),
 \label{eq:pure21-temporal-selection}\\
 (0,2k_0)+(\omega_0,-k_0)&=(\omega_0,k_0),\notag\\
 (0,-2k_0)+(\omega_0,k_0)&=(\omega_0,-k_0).
 \label{eq:pure21-spatial-selection}
\end{align}
Projection by the adjoint critical eigenvector therefore gives
\begin{align}
 \eta_t&=2\Psi^*\cQ(\overline\Phi,p_2),\notag\\
 \eta_x^+&=2\Psi^*\cQ(\Phi,s_2^+),\qquad
 \eta_x^-=2\Psi^*\cQ(\Phi,s_2^-).
 \label{eq:pure21-eta}
\end{align}
Reflection symmetry and a choice of temporal and spatial origins reduce
these to $\eta_t,\eta_x\geq0$ in \eqref{eq:forced}.

The symbol $f$ was not a coefficient of the original unforced amplitude
system.  It was introduced only as shorthand for a possible
amplitude-independent critical residual in a more general forced reduction.
If such a third-order residual had critical coefficients $R_+$ and $R_-$,
the solvability projection would give
\begin{equation}
 f_+=\Psi^*R_+,
 \qquad
 f_-=\Psi^*R_-.
 \label{eq:pure21-f-origin}
\end{equation}
The calculation below shows that $R_+=R_-=0$ for the present exact
$2{:}1$-or-absent inputs, so neither $f_+$ nor $f_-$ appears in
\eqref{eq:forced}.

There are two separate reasons that no additive term occurs. First, a
temporal forcing harmonic has pair $(2\omega_0,0)$ and a spatial forcing
harmonic has pair $(0,\pm2k_0)$; neither equals either critical pair
$(\omega_0,\pm k_0)$. Second, their first mixed product has order
\begin{equation}
 \epsilon^n\epsilon^m=\epsilon^4
 \label{eq:pure21-order-four}
\end{equation}
and pair $(2\omega_0,\pm2k_0)$, so it is neither third order nor critical.
More strongly, every amplitude-independent monomial formed only from the
doubled forcing harmonics has pair
\begin{equation}
 (2j\omega_0,2\ell k_0),\qquad j,\ell\in\mathbb Z,
 \label{eq:pure21-even-lattice}
\end{equation}
which cannot equal $(\omega_0,\pm k_0)$. Hence an analytic nonlinearity
cannot generate an additive critical term from exact $2{:}1$ separated
forcing at any order. In this report's regime,
\begin{equation}
 \boxed{f=0.}
 \label{eq:pure21-fzero}
\end{equation}
An additive term would require a $1{:}1$ spatio-temporal harmonic, leakage
outside the doubled forcing lattice, or explicit translation-breaking
heterogeneity; none is assumed here.

If the temporal frequency is
$2\omega_0+\epsilon^2\delta_t$, the parametric terms initially contain
$\ee^{\ii\delta_tT}$. The transformation
\begin{equation}
 A=\widetilde A\ee^{\ii\delta_tT/2},\qquad
 B=\widetilde B\ee^{\ii\delta_tT/2}
 \label{eq:pure21-detuning-transform}
\end{equation}
makes the system autonomous and replaces $a$ by
$a-\ii\delta_t/2$. We therefore allow $a,b,c\in\mathbb C$ below and use
\eqref{eq:forced} without tildes.

\subsection{Exact polar equations and the definition of locking}

Write
\begin{equation}
 A=r\ee^{\ii\phi},\qquad B=\rho\ee^{\ii\psi},
 \qquad r,\rho\geq0,
 \label{eq:general-polar-form}
\end{equation}
and separate
\begin{align}
 a+br^2+c\rho^2&=\alpha_A+\ii\beta_A,&
 a+b\rho^2+cr^2&=\alpha_B+\ii\beta_B.
 \label{eq:polar-local-coefficients}
\end{align}
For $r,\rho>0$, multiplication of \eqref{eq:forced} by
$\ee^{-\ii\phi}$ and $\ee^{-\ii\psi}$ gives exactly
\begin{align}
 \dot r={}&r\alpha_A+\eta_t\rho\cos(\phi+\psi)
 +\eta_x\rho\cos(\phi-\psi),
 \label{eq:polar-r}\\
 r\dot\phi={}&r\beta_A-\eta_t\rho\sin(\phi+\psi)
 -\eta_x\rho\sin(\phi-\psi),
 \label{eq:polar-phi}\\
 \dot\rho={}&\rho\alpha_B+\eta_t r\cos(\phi+\psi)
 +\eta_xr\cos(\phi-\psi),
 \label{eq:polar-rho}\\
 \rho\dot\psi={}&\rho\beta_B-\eta_t r\sin(\phi+\psi)
 +\eta_xr\sin(\phi-\psi).
 \label{eq:polar-psi}
\end{align}
A phase-locked equilibrium satisfies
\begin{equation}
 \dot r=\dot\rho=\dot\phi=\dot\psi=0.
 \label{eq:general-locking-definition}
\end{equation}
The physical carrier still oscillates at frequency $\omega_0$; locking means
that both slow amplitudes and both phase offsets relative to the forcing are
constant.

Put
\begin{equation}
 \Theta=\phi+\psi,\qquad \Delta=\phi-\psi.
 \label{eq:phase-sum-difference}
\end{equation}
The radial equations contain
\begin{equation}
 C=\eta_t\cos\Theta+\eta_x\cos\Delta
 \label{eq:phase-C}
\end{equation}
and are equivalent to
\begin{equation}
 r\alpha_A+\rho C=0,\qquad
 \rho\alpha_B+rC=0.
 \label{eq:general-radial-locking}
\end{equation}
Consequently
\begin{equation}
 C=-\frac{r}{\rho}\alpha_A
   =-\frac{\rho}{r}\alpha_B,\qquad
 r^2\alpha_A=\rho^2\alpha_B.
 \label{eq:general-radial-compatibility}
\end{equation}
The phase equations give
\begin{align}
 T:=\eta_t\sin\Theta
 &=\frac12\left(\frac{r}{\rho}\beta_A
                 +\frac{\rho}{r}\beta_B\right),
 \label{eq:phase-T}\\
 X:=\eta_x\sin\Delta
 &=\frac12\left(\frac{r}{\rho}\beta_A
                 -\frac{\rho}{r}\beta_B\right).
 \label{eq:phase-X}
\end{align}
For each $\varepsilon_t,\varepsilon_x\in\{-1,1\}$, positive amplitudes give
a locked state if and only if
\begin{align}
 |T|&\leq\eta_t,\qquad |X|\leq\eta_x,
 \label{eq:phase-sine-bounds}\\
 -\frac{r}{\rho}\alpha_A
 &=\varepsilon_t\sqrt{\eta_t^2-T^2}
  +\varepsilon_x\sqrt{\eta_x^2-X^2},
 \label{eq:general-phase-amplitude-closure}
\end{align}
together with the last equality in
\eqref{eq:general-radial-compatibility}. The phases are then
\begin{align}
 \Theta&=\operatorname{atan2}\!\left(
 T,\varepsilon_t\sqrt{\eta_t^2-T^2}\right),\notag\\
 \Delta&=\operatorname{atan2}\!\left(
 X,\varepsilon_x\sqrt{\eta_x^2-X^2}\right),
 \label{eq:general-phase-reconstruction}\\
 \phi&=\frac{\Theta+\Delta}{2},\qquad
 \psi=\frac{\Theta-\Delta}{2}.
 \label{eq:individual-phase-reconstruction}
\end{align}
These conditions are necessary and sufficient, not merely necessary.
If $\eta_t=0$, the first sine condition is replaced by $T=0$, the term
$\eta_t\cos\Theta$ is zero, and $\Theta$ is a neutral phase; the first
$\operatorname{atan2}$ formula is then omitted.  If $\eta_x=0$, the second
sine condition is replaced by $X=0$, the term $\eta_x\cos\Delta$ is zero,
and $\Delta$ is neutral; the second $\operatorname{atan2}$ formula is
omitted.  After assigning arbitrary values to the neutral phases,
\eqref{eq:individual-phase-reconstruction} still reconstructs $\phi$ and
$\psi$.  Thus no undefined $\operatorname{atan2}(0,0)$ is used when an input
is absent.

\subsection{A factorization that exhausts the positive amplitudes}

Let $a_r=\RePart a$, $b_r=\RePart b$, and $c_r=\RePart c$. Expanding the
last condition in \eqref{eq:general-radial-compatibility} gives
\begin{align}
 0={}&r^2(a_r+b_rr^2+c_r\rho^2)
 -\rho^2(a_r+b_r\rho^2+c_rr^2)\notag\\
 ={}&(r^2-\rho^2)\left[a_r+b_r(r^2+\rho^2)\right].
 \label{eq:forced-amplitude-factorization}
\end{align}
Therefore every positive equilibrium satisfies
\begin{equation}
 \boxed{r=\rho}
\end{equation}
or, when $b_r\ne0$, it satisfies
\begin{equation}
 \boxed{r^2+\rho^2=-\frac{a_r}{b_r}}.
 \label{eq:equal-or-unequal}
\end{equation}
If $b_r=0$ and $a_r\ne0$, no unequal-amplitude state exists. If
$a_r=b_r=0$, the factorization is degenerate and the necessary-and-sufficient
closure \eqref{eq:phase-sine-bounds}--
\eqref{eq:general-phase-amplitude-closure} is used directly.

For the unequal class define only
\begin{equation}
 R=r^2+\rho^2=-\frac{a_r}{b_r},\qquad
 \Pi=r\rho,\qquad
 D=r^2-\rho^2=\pm\sqrt{R^2-4\Pi^2},
 \label{eq:unequal-invariants}
\end{equation}
where $R>0$ and $0<\Pi<R/2$. The quantities in the phase closure reduce to
\begin{align}
 C&=(b_r-c_r)\Pi,
 \label{eq:unequal-C}\\
 T(\Pi)&=\frac{a_iR+b_iR^2+2(c_i-b_i)\Pi^2}{2\Pi},
 \label{eq:unequal-T}\\
 X(\Pi,D)&=\frac{D(a_i+b_iR)}{2\Pi},
 \label{eq:unequal-X}
\end{align}
where $a_i=\ImPart a$, $b_i=\ImPart b$, and $c_i=\ImPart c$. Hence all
unequal-amplitude equilibria are the roots $0<\Pi<R/2$ of
\begin{equation}
 (b_r-c_r)\Pi
 =\varepsilon_t\sqrt{\eta_t^2-T(\Pi)^2}
 +\varepsilon_x\sqrt{\eta_x^2-X(\Pi,D)^2},
 \label{eq:unequal-scalar-equation}
\end{equation}
subject to $|T(\Pi)|\leq\eta_t$ and
$|X(\Pi,D)|\leq\eta_x$. The two signs of $D$ give the exchanged states.
Then
\begin{equation}
 r^2=\frac{R+D}{2},\qquad \rho^2=\frac{R-D}{2},
 \label{eq:unequal-reconstruction}
\end{equation}
and the phases follow from
\eqref{eq:general-phase-reconstruction}--
\eqref{eq:individual-phase-reconstruction}.

The remaining scalar problem can be reduced to a quartic without losing the
sign information.  Put $z=\Pi^2$, $d_r=b_r-c_r$, and
\begin{equation}
 L=a_i+b_iR,
 \qquad
 N(z)=a_iR+b_iR^2+2(c_i-b_i)z.
 \label{eq:unequal-polynomial-shortcuts}
\end{equation}
The two squared radical arguments, after multiplication by $4z$, are
\begin{align}
 U_0(z)&=4\eta_t^2z-N(z)^2,
 \label{eq:unequal-U0}\\
 V_0(z)&=4\eta_x^2z-(R^2-4z)L^2.
 \label{eq:unequal-V0}
\end{align}
Thus the unsquared equation has
\begin{equation}
 C=d_r\sqrt z,
 \qquad
 \eta_t^2-T^2=\frac{U_0(z)}{4z},
 \qquad
 \eta_x^2-X^2=\frac{V_0(z)}{4z}.
 \label{eq:unequal-radical-arguments}
\end{equation}
Eliminating the two square roots gives the explicit quartic equation
\begin{equation}
 16d_r^2z^2V_0(z)
 -\left[4d_r^2z^2+V_0(z)-U_0(z)\right]^2=0.
 \label{eq:unequal-quartic}
\end{equation}
Every unequal equilibrium gives a root of \eqref{eq:unequal-quartic}.  A
quartic root is retained exactly when
\begin{equation}
 0<z<\frac{R^2}{4},
 \qquad U_0(z)\geq0,
 \qquad V_0(z)\geq0,
 \label{eq:unequal-quartic-filter}
\end{equation}
and it satisfies the original signed equation
\eqref{eq:unequal-scalar-equation} for at least one pair
$(\varepsilon_t,\varepsilon_x)$.  This final unsquared test removes every
extraneous root introduced by squaring.  Consequently there are at most four
admissible unequal products $\Pi$ at generic cubic order, and the two signs
of $D$ reconstruct the exchanged states.

\subsection{Closed equal-amplitude solutions at every phase}

Let $r=\rho>0$ and first assume $\eta_x>0$. Then
\eqref{eq:phase-X} gives $\sin\Delta=0$, so
\begin{equation}
 B=\varepsilon A,\qquad \varepsilon\in\{1,-1\},\qquad
 A=r\ee^{\ii\phi},\qquad x=r^2.
 \label{eq:equal-general-phase-ansatz}
\end{equation}
With $s=b+c$ and
\begin{equation}
 h_\varepsilon(x)=a+\varepsilon\eta_x+sx,
 \label{eq:equal-h}
\end{equation}
the equilibrium equation is
\begin{equation}
 h_\varepsilon(x)+\varepsilon\eta_t\ee^{-2\ii\phi}=0.
 \label{eq:equal-general-phase-balance}
\end{equation}
For $\eta_t>0$, taking the modulus gives
\begin{equation}
 |s|^2x^2+2\RePart\!\left[
 (a+\varepsilon\eta_x)\overline s\right]x
 +|a+\varepsilon\eta_x|^2-\eta_t^2=0.
 \label{eq:equal-general-phase-quadratic}
\end{equation}
Put
\begin{equation}
 D_\varepsilon=
 \RePart\!\left[(a+\varepsilon\eta_x)\overline s\right]^2
 -|s|^2\left(|a+\varepsilon\eta_x|^2-\eta_t^2\right).
 \label{eq:equal-general-phase-discriminant}
\end{equation}
For $s\ne0$, all squared amplitudes are
\begin{equation}
 x_{\varepsilon,\pm}=
 \frac{-\RePart[(a+\varepsilon\eta_x)\overline s]
 \pm\sqrt{D_\varepsilon}}{|s|^2},
 \label{eq:equal-general-phase-roots}
\end{equation}
retaining exactly the roots with $D_\varepsilon\geq0$ and
$x_{\varepsilon,\pm}>0$. The corresponding phases are
\begin{align}
 \ee^{-2\ii\phi}
 &=-\frac{h_\varepsilon(x)}{\varepsilon\eta_t},
 \label{eq:equal-general-phase-exponential}\\
 \phi&=-\frac12\arg\!\left[
 -\frac{h_\varepsilon(x)}{\varepsilon\eta_t}\right]\pmod{\pi}.
 \label{eq:equal-general-phase-angle}
\end{align}
The two values modulo $2\pi$ differ by $\pi$ and give the sign-related
states $(A,B)$ and $(-A,-B)$.

If $s=0$, the modulus equation is independent of $x$.  There is no positive
equal-amplitude branch unless
\begin{equation}
 |a+\varepsilon\eta_x|=\eta_t.
 \label{eq:equal-degenerate-condition}
\end{equation}
When this equality holds and $\eta_t>0$, every $x>0$ satisfies the
cubic-order modulus condition and its phase is given by
\eqref{eq:equal-general-phase-exponential}.  If $\eta_t=0$, the undivided
condition is $a+\varepsilon\eta_x=0$ and the common phase is neutral.  Each
continuum is a degeneracy of the cubic truncation and must be resolved by
quintic terms.

For temporal forcing alone, $\eta_x=0$ and $\eta_t>0$, the positive
equal-amplitude family
does not require $B=\pm A$.  Instead, for $x=r^2$ it is
\begin{equation}
 A=r\ee^{\ii(\Theta+\Delta)/2},\qquad
 B=r\ee^{\ii(\Theta-\Delta)/2},\qquad
 |a+sx|=\eta_t,\qquad
 \ee^{-\ii\Theta}=-\frac{a+sx}{\eta_t},
 \label{eq:temporal-only-equal-family}
\end{equation}
where $\Delta$ is arbitrary.  This gives the general locked phase rather
than only the special choices $B=\pm A$.

The cases in which one or both inputs are absent follow without a limiting
ambiguity:
\begin{center}
\begin{tabular}{@{}L{0.18\textwidth}L{0.36\textwidth}L{0.34\textwidth}@{}}
\toprule
Forcing present & Equal-amplitude condition & Phase information\\
\midrule
Temporal and spatial &
$|a+\varepsilon\eta_x+sx|=\eta_t$ &
$\Delta=0$ or $\pi$; $\phi$ from
\eqref{eq:equal-general-phase-angle}\\
Temporal only &
$|a+sx|=\eta_t$ &
$\Theta$ is locked by
$a+sx+\eta_t\ee^{-\ii\Theta}=0$; $\Delta$ is neutral\\
Spatial only &
$a+\varepsilon\eta_x+sx=0$ &
$\Delta=0$ or $\pi$; the common phase is neutral\\
Neither &
$a+sx=0$ for a stationary complex equilibrium &
both group phases are neutral; Section~\ref{sec:equilibria} gives the more
general relative equilibria\\
\bottomrule
\end{tabular}
\end{center}
With no forcing, a radial equilibrium whose imaginary frequency correction
is nonzero is a relative equilibrium rather than a stationary point in the
chosen rotating frame.

\subsection{Exceptional one-sided states}

Polar coordinates omit $r=0$ or $\rho=0$, so return to
\eqref{eq:forced}. If $A=0$ and $B=\rho\ee^{\ii\psi}\ne0$, then
\begin{equation}
 \eta_t\overline B+\eta_xB=0,\qquad
 a+b\rho^2=0.
 \label{eq:one-sided-conditions}
\end{equation}
If at least one parametric forcing is present, the first equation implies
\begin{equation}
 \eta_t=\eta_x>0,\qquad
 \ee^{2\ii\psi}=-1,
 \label{eq:one-sided-phase}
\end{equation}
and, when $b\ne0$, the second requires the complex number $-a/b$ to be real
and positive, with $\rho^2=-a/b$.  If $b=0$, it instead requires $a=0$ and
leaves every $\rho>0$ unresolved at cubic order.  Exchanging $A$ and $B$
gives the other one-sided state.
If $\eta_t\ne\eta_x$, no nonzero one-sided equilibrium exists. These
codimension conditions explain why generic parametric forcing converts
unforced traveling waves into two-component locked states.
When $\eta_t=\eta_x=0$, the first condition in
\eqref{eq:one-sided-conditions} disappears and the usual unforced
one-sided traveling-wave relative equilibria of
Section~\ref{sec:equilibria} return.

\subsection{Why the equilibrium list is exhaustive}

Every point $(A,B)\in\mathbb C^2$ belongs to exactly one of the disjoint
amplitude classes
\begin{equation}
 (|A|,|B|)=(0,0),\qquad
 |A||B|=0<|A|+|B|,\qquad
 |A|=|B|>0,\qquad
 |A|\ne|B|,\quad |A||B|>0.
 \label{eq:forced-exhaustive-partition}
\end{equation}
The first class is rest.  The second is exactly the one-sided calculation
\eqref{eq:one-sided-conditions}.  In the third class, the phase equations
give either \eqref{eq:equal-general-phase-roots}--
\eqref{eq:equal-general-phase-angle}, the temporal-only family
\eqref{eq:temporal-only-equal-family}, or one of the absent-input rows in
the table above.  In the fourth class,
\eqref{eq:forced-amplitude-factorization} forces
\eqref{eq:equal-or-unequal}, and the necessary-and-sufficient phase closure
reduces it to \eqref{eq:unequal-scalar-equation}.  Hence these formulas omit
no complex equilibrium.  The exceptional coefficient cases $b_r=0$ or
$a_r=b_r=0$ are not divided by; they are decided directly by
\eqref{eq:general-radial-compatibility} and
\eqref{eq:general-phase-amplitude-closure}.

\subsection{Closed real-coefficient classification}

For real $a,b,c$ and real $A,B$, put
\begin{equation}
 q=\eta_t+\eta_x,\qquad s=b+c,\qquad d=b-c.
 \label{eq:forcing-combinations}
\end{equation}
Adding and subtracting the two equilibrium equations gives
\begin{align}
 0&=(A-B)\left[a-q+b(A^2+B^2)+dAB\right],
 \label{eq:forced-difference}\\
 0&=(A+B)\left[a+q+b(A^2+B^2)-dAB\right].
 \label{eq:forced-sum}
\end{align}
Therefore the complete real list is:
\begin{align}
 &(A,B)=(0,0),\notag\\
 &(A,B)=(r,r),\qquad r^2=-\frac{a+q}{s}>0,
 \label{eq:real-in-phase}\\
 &(A,B)=(r,-r),\qquad r^2=-\frac{a-q}{s}>0,
 \label{eq:real-antiphase}
\end{align}
plus exchange-broken states. When $bd\ne0$, the latter satisfy
\begin{equation}
 A^2+B^2=-\frac{a}{b},\qquad AB=\frac{q}{d}.
 \label{eq:real-broken-invariants}
\end{equation}
Writing $U=-a/b$ and $P=q/d$, they exist exactly when
\begin{equation}
 U+2P\geq0,\qquad U-2P\geq0,
 \label{eq:real-broken-existence}
\end{equation}
and all sign-related and exchanged roots are
\begin{equation}
 A=\frac{\sigma\sqrt{U+2P}+\tau\sqrt{U-2P}}{2},\qquad
 B=\frac{\sigma\sqrt{U+2P}-\tau\sqrt{U-2P}}{2},
 \quad \sigma,\tau\in\{-1,1\}.
 \label{eq:real-broken-roots}
\end{equation}
Equalities in \eqref{eq:real-broken-existence} merge this family with an
equal or anti-equal branch. If $b=0$, $d=0$, or $s=0$, the undivided
factorization \eqref{eq:forced-difference}--\eqref{eq:forced-sum} gives the
corresponding degenerate cases without losing a branch.

\subsection{Exact stability: rest and every equal-amplitude branch}

At the rest state, exchange-symmetric and exchange-antisymmetric
perturbations satisfy
\begin{equation}
 \dot z_\kappa=(a+\kappa\eta_x)z_\kappa
 +\kappa\eta_t\overline z_\kappa,\qquad \kappa\in\{1,-1\}.
 \label{eq:rest-exchange-modes}
\end{equation}
Their characteristic polynomials and eigenvalues are
\begin{align}
 0={}&\lambda^2-2\RePart(a+\kappa\eta_x)\lambda
 +|a+\kappa\eta_x|^2-\eta_t^2,
 \label{eq:rest-characteristic}\\
 \lambda_{\kappa,\pm}
 ={}&\RePart a+\kappa\eta_x
 \pm\sqrt{\eta_t^2-(\ImPart a)^2}.
 \label{eq:rest-eigenvalues}
\end{align}
The square root is imaginary when its argument is negative. The rest state
is asymptotically stable if and only if, for both $\kappa=\pm1$,
\begin{equation}
 \RePart(a+\kappa\eta_x)<0,\qquad
 |a+\kappa\eta_x|^2>\eta_t^2.
 \label{eq:rest-stability}
\end{equation}

Now let $B=\varepsilon A$, $A=r\ee^{\ii\phi}$, and $x=r^2$. The full
$4\times4$ linearization splits exactly into the two exchange sectors
$\delta B=\kappa\delta A$. In sector $\kappa$ it is the scalar
conjugate-linear equation
\begin{equation}
 \delta\dot A=P_\kappa\delta A+Q_\kappa\delta\overline A,
 \label{eq:equal-sector-linearization}
\end{equation}
where
\begin{align}
 P_\kappa&=
 a+\left[2b+c+\kappa\varepsilon c\right]x+\kappa\eta_x,
 \label{eq:equal-sector-P}\\
 Q_\kappa&=
 x\ee^{2\ii\phi}(b+\kappa\varepsilon c)+\kappa\eta_t.
 \label{eq:equal-sector-Q}
\end{align}
The two eigenvalues in that sector obey
\begin{equation}
 \lambda^2-2\RePart P_\kappa\,\lambda
 +|P_\kappa|^2-|Q_\kappa|^2=0
 \label{eq:equal-sector-characteristic}
\end{equation}
and therefore
\begin{equation}
 \lambda_{\kappa,\pm}
 =\RePart P_\kappa
 \pm\sqrt{|Q_\kappa|^2-(\ImPart P_\kappa)^2}.
 \label{eq:equal-sector-eigenvalues}
\end{equation}
Thus an equal-amplitude locked state is asymptotically stable if and only if,
for each $\kappa=\pm1$,
\begin{equation}
 \boxed{\RePart P_\kappa<0,\qquad
 |P_\kappa|^2>|Q_\kappa|^2.}
 \label{eq:equal-sector-stability}
\end{equation}
If one forcing is absent, the continuous symmetry it leaves unbroken
produces the expected zero eigenvalue; stability is then orbital after
removing that neutral group direction.

\subsection{Exact stability at every remaining complex root}

Let the right-hand side of the first equation in \eqref{eq:forced} be
$F_A$. Its Wirtinger derivatives are
\begin{align}
 \partial_A F_A&=a+2b|A|^2+c|B|^2,
 &\partial_{\overline A}F_A&=bA^2,\notag\\
 \partial_B F_A&=cA\overline B+\eta_x,
 &\partial_{\overline B}F_A&=cAB+\eta_t.
 \label{eq:wirtinger}
\end{align}
The derivatives of $F_B$ follow by exchanging $A$ and $B$. The real matrix
of $z\mapsto Pz+Q\overline z$ is
\begin{equation}
 \mathcal R(P,Q)=
 \begin{pmatrix}
 \RePart(P+Q)&-\ImPart(P-Q)\\
 \ImPart(P+Q)& \RePart(P-Q)
 \end{pmatrix}.
 \label{eq:realification}
\end{equation}
Using \eqref{eq:realification} on the four derivative pairs in
\eqref{eq:wirtinger} and its exchanged copy produces the exact real
$4\times4$ Jacobian
\begin{equation}
 J=\begin{pmatrix}
 \mathcal R(a+2b|A|^2+c|B|^2,bA^2)&
 \mathcal R(cA\overline B+\eta_x,cAB+\eta_t)\\
 \mathcal R(cB\overline A+\eta_x,cAB+\eta_t)&
 \mathcal R(a+2b|B|^2+c|A|^2,bB^2)
 \end{pmatrix}.
 \label{eq:full-forced-Jacobian}
\end{equation}
This formula applies at any root reconstructed above, including the
codimension-one-sided states.

For a coefficient-only stability test, write
\begin{equation}
 \det(\lambda I-J)=
 \lambda^4+h_1\lambda^3+h_2\lambda^2+h_3\lambda+h_4,
 \label{eq:quartic-characteristic}
\end{equation}
where no eigenvalue calculation is required:
\begin{align}
 h_1&=-\operatorname{tr}J,\notag\\
 h_2&=\frac12\left[(\operatorname{tr}J)^2-\operatorname{tr}(J^2)\right],
 \notag\\
 h_3&=-\frac16\left[(\operatorname{tr}J)^3
 -3\operatorname{tr}J\,\operatorname{tr}(J^2)
 +2\operatorname{tr}(J^3)\right],\notag\\
 h_4&=\det J.
 \label{eq:quartic-coefficients}
\end{align}
The root is asymptotically stable if and only if the quartic
Routh--Hurwitz inequalities hold:
\begin{equation}
 h_1,h_2,h_3,h_4>0,\qquad
 h_1h_2>h_3,\qquad
 h_1h_2h_3>h_3^2+h_1^2h_4.
 \label{eq:quartic-hurwitz}
\end{equation}
Equations \eqref{eq:unequal-scalar-equation} and
\eqref{eq:quartic-hurwitz} therefore give a complete existence-and-stability
algorithm for every unequal complex phase-locked state.  Applying the same
Jacobian and Hurwitz test to \eqref{eq:one-sided-conditions} gives the exact
stability of every exceptional one-sided state.

\subsection{Application to the computed three- and five-population coefficients}

For the real-coefficient projection, take $a=0.01$ and
$\eta_t=\eta_x=0.04$. The real parts of
\eqref{eq:numerical-coefficients} give
\begin{equation}
 (A,B)=(0,0),\qquad (A,B)=(r,r),\qquad(A,B)=(-r,-r),
 \qquad r=0.23173.
\end{equation}
The rest state is unstable. Each nonzero state has eigenvalues
\begin{equation}
 (-0.18000,-0.00412,-0.08000,-0.08000)
\end{equation}
and is asymptotically stable within this real-coefficient projection. These
real states are not equilibria of the full complex-coefficient normal form.

At zero detuning, retain the complete complex values of $b,c$ and apply
\eqref{eq:equal-general-phase-roots}--
\eqref{eq:equal-general-phase-angle}. The exact sector formula
\eqref{eq:equal-sector-eigenvalues} gives:
\begin{center}
\begin{tabular}{@{}lrrrrl@{}}
\toprule
Model & $\varepsilon$ & $r$ & $\phi$ (degrees) & $\max\RePart\lambda$ & Stability\\
\midrule
Three-population & $+1$ & 0.20564 & 150.732 & 0.00100 & unstable\\
Three-population & $+1$ & 0.07843 &  93.563 & 0.07839 & unstable\\
Three-population & $-1$ & 0.07617 &  86.639 & 0.08000 & unstable\\
Five-population  & $+1$ & 0.14949 & 172.622 & $-0.07861$ & stable\\
Five-population  & $+1$ & 0.05024 &  90.824 & 0.07991 & unstable\\
Five-population  & $-1$ & 0.05016 &  89.178 & 0.08000 & unstable\\
\bottomrule
\end{tabular}
\end{center}
Thus all equal-amplitude branches of the full three-population model are
unstable at this forcing, although the large branch is only weakly unstable.
The five-population model has the stable complex locked state
\begin{equation}
 (A,B)=\left(0.14949\ee^{\ii\,172.622^\circ},
              0.14949\ee^{\ii\,172.622^\circ}\right)
 \label{eq:five-complex-locked-state}
\end{equation}
and its sign-related partner.

The equal-amplitude table is not the complete numerical list.  For both
models the one-sided condition fails because $-a/b$ is not real.  Solving
the quartic \eqref{eq:unequal-quartic} and applying the unsquared filter
\eqref{eq:unequal-scalar-equation} gives the following representatives:
\begin{center}
\begin{tabular}{@{}lrrrrrl@{}}
\toprule
Model & $r$ & $\rho$ & $\phi$ (degrees) & $\psi$ (degrees) &
$\max\RePart\lambda$ & Stability\\
\midrule
Three-population & 0.06978 & 0.29011 & $-24.559$ & $-66.184$ &
 0.01390 & unstable\\
Three-population & 0.15984 & 0.25197 & $-26.250$ & $-35.442$ &
 $-0.00181$ & stable\\
Five-population & 0.00077 & 0.09520 & $-16.561$ & 89.904 &
 0.07233 & unstable\\
\bottomrule
\end{tabular}
\end{center}
Each row has an exchanged partner obtained by interchanging
$(r,\phi)$ with $(\rho,\psi)$, and each of those has a sign-related partner
obtained by adding $180^\circ$ to both phases.  Substitution into the full
four-real-dimensional system gives residuals below $2\times10^{-17}$ before
rounding the displayed values.

Therefore the complete calculation changes the activation-pattern
conclusion for the three-population coefficients: although every
equal-amplitude locked branch is unstable, there is a stable
exchange-broken phase-locked pair.  For the five-population coefficients,
the stable state remains the equal-amplitude branch
\eqref{eq:five-complex-locked-state}; its unequal pair is unstable.

\begin{figure}[ht]
\centering
\includegraphics[width=\textwidth]{output/simulations/amplitude_phase_diagrams.png}
\caption{Amplitude phase diagrams using the computed real cubic coefficients.
White circles are stable states and the cross is the unstable unforced
standing wave.}
\label{fig:phase}
\end{figure}
